\documentclass[a4paper]{article}

\usepackage{graphicx}
\usepackage{amsmath,amssymb}
\usepackage{minted}
\usepackage{multirow}
\usepackage{float}
\usepackage{algorithm}
\usepackage{algpseudocode}
\usepackage{todonotes}
\usepackage{forest}
\usepackage{xstring}
\usepackage{xcolor}
\usepackage[most]{tcolorbox}
\usepackage{xparse}
\usepackage{natbib}

\usetikzlibrary{graphs,graphdrawing}
\usegdlibrary{layered}

% for external graphviz
\input{/home/donkarlo/Dropbox/repo/nd_language_ai_project/src/nd_language_ai/formal/computer/programming/tex/latex/code_clips/external_graphviz.tex}

% for tags
\input{/home/donkarlo/Dropbox/repo/nd_language_ai_project/src/nd_language_ai/formal/computer/programming/tex/latex/part/control_sequence/kind/semtag/semtag.tex}

% for links
\input{/home/donkarlo/Dropbox/repo/nd_language_ai_project/src/nd_language_ai/formal/computer/programming/tex/latex/code_clips/seehere.tex}

\title{Source and Visual typing latency and Source scroll stability fix}
\date{2026-09-17}
\author{Mohammad Rahmani}

\begin{document}
   \maketitle

   \section{Problems}
      Source typing could pause when the graphical formula/table preview controller
      launched LuaLaTeX synchronously on the GUI thread. The controller also removed
      all preview bottom margins before rebuilding them. Near the bottom of a Source
      document this temporarily reduced the scrollbar maximum, so Qt clamped the
      current scroll position and the page visibly jumped upward when the margins
      were added again.

      Visual edits were also propagated aggressively enough that full Visual
      serialization and Source synchronization could start during ordinary typing
      pauses.

   \section{Fix}
      \begin{itemize}
         \item Formula/table Source preview rendering now runs in a dedicated
         short-lived worker thread. The GUI thread only extracts spans, maintains
         cached images, and applies the final overlay layout.
         \item While a render is running, newer edits replace the pending request.
         Stale render results are cached when useful but are never committed as the
         active Source layout.
         \item Existing Source preview margins remain active until the replacement
         layout is ready. Margin updates are differential instead of clearing and
         rescanning the entire document.
         \item The Source controller preserves the visible top anchor and explicitly
         keeps the scrollbar pinned to the bottom when the user was already at the
         bottom. This removes the end-of-file upward jump reproduced in
         \texttt{neural\_cognitive\_architecture.tex}.
         \item The normal Source content debounce is increased from 70 ms to 240 ms.
         \item The normal Visual propagation debounce is increased from 180 ms to
         650 ms, and the large Visual debounce from 650 ms to 1200 ms.
         \item The PDF preview debounce is increased from 120 ms to 260 ms so preview
         work starts after a short typing idle period instead of competing with
         keystrokes.
      \end{itemize}

   \section{Data safety}
      The canonical LaTeX source remains unchanged by the preview controller. The
      worker renders temporary standalone snippets only. Source/Visual mode changes
      still use the existing synchronous pending-edit flush boundary, so the larger
      debounce values do not discard final Visual keystrokes.

   \bibliographystyle{plainnat}
   \bibliography{/home/donkarlo/Dropbox/repo/nd_shared_research/bibliography/bibliography.bib}
\end{document}
